\documentclass[11pt, a4paper]{article}

% ── Packages ──────────────────────────────────────────────────────────────────
\usepackage{amsmath, amssymb, amsthm}
\usepackage{booktabs}
\usepackage{enumitem}
\usepackage[T1]{fontenc}
\usepackage{geometry}
\usepackage[utf8]{inputenc}
\usepackage{microtype}
\usepackage{parskip}
\usepackage{titlesec}
\usepackage{titletoc}
\usepackage{xcolor}

% hyperref last
\usepackage[
  colorlinks = true,
  linkcolor  = teal!70!black,
  citecolor  = teal!70!black,
  urlcolor   = blue!60!black,
  pdftitle   = {SierpSphere Grammar Notation},
]{hyperref}

% ── Page geometry ─────────────────────────────────────────────────────────────
\geometry{margin=2.8cm, top=3.2cm, bottom=3.2cm}

% ── Colours ───────────────────────────────────────────────────────────────────
\definecolor{accent}{RGB}{0, 120, 120}

% ── Section formatting ────────────────────────────────────────────────────────
\titleformat{\section}
  {\large\bfseries\color{accent}}{\thesection.}{0.6em}{}[\color{accent}\titlerule]
\titleformat{\subsection}
  {\normalsize\bfseries\color{accent!80!black}}{\thesubsection}{0.5em}{}
\titleformat{\paragraph}[runin]
  {\normalsize\bfseries}{}{0em}{}[.\enspace]

\titlespacing{\section}{0pt}{1.8em}{0.6em}
\titlespacing{\subsection}{0pt}{1.2em}{0.3em}

% ── TOC ───────────────────────────────────────────────────────────────────────
\titlecontents{section}[1.5em]
  {\smallskip\bfseries}
  {\contentslabel{1.5em}}{\hspace*{-1.5em}}
  {\titlerule*[6pt]{.}\contentspage}
\titlecontents{subsection}[3.5em]
  {}{\contentslabel{2em}}{}
  {\titlerule*[6pt]{.}\contentspage}

% ── Grammar macros ────────────────────────────────────────────────────────────
\newcommand{\G}[1]{\mathsf{#1}}
\newcommand{\prim}[1]{\mathrm{#1}}
\newcommand{\sub}{\ominus}
\newcommand{\add}{\oplus}
\newcommand{\intr}{\cap}
\newcommand{\half}{\tfrac{1}{2}}
\newcommand{\code}[1]{\texttt{#1}}

% ── Title ─────────────────────────────────────────────────────────────────────
\title{%
  {\color{accent}\rule{\linewidth}{2pt}}\\[6pt]
  \textbf{\LARGE SierpSphere Grammar Notation}\\[4pt]
  \large A compact symbolic language for fractal SDF grammars\\[6pt]
  {\color{accent}\rule{\linewidth}{1pt}}
}
\author{\href{https://github.com/carlok/sierpsphere}%
        {\texttt{github.com/carlok/sierpsphere}}}
\date{2026}

% ─────────────────────────────────────────────────────────────────────────────
\begin{document}

\maketitle
\thispagestyle{empty}

\vspace{0.5em}
{\small\tableofcontents}
\vspace{1em}

% ══ 1 ══════════════════════════════════════════════════════════════════════════
\section{Motivation}
\label{sec:motivation}

A SierpSphere grammar is fully specified by three things: a \emph{symmetry
group}, a \emph{seed primitive}, and an ordered list of \emph{iteration steps}.
The JSON representation is verbose; this notation condenses a grammar to a
single line amenable to comparison, mutation logging, and typeset publication.
For the broader context see
\href{journey.pdf}{\emph{Growing Metal Lace}}.

% ══ 2 ══════════════════════════════════════════════════════════════════════════
\section{Alphabet}
\label{sec:alphabet}

\subsection{Symmetry groups}
\label{subsec:groups}

Symmetry groups are denoted by their Schoenflies symbols:

\[
  \G{T_d} \quad \G{O_h} \quad \G{I_h}
\]

corresponding to the \textit{tetrahedral} (4 axes),
\textit{octahedral} (6 axes), and \textit{icosahedral} (12 axes) point groups
respectively.

\subsection{Primitives}
\label{subsec:primitives}

\[
  \prim{S} = \text{sphere} \qquad
  \prim{C} = \text{cube}   \qquad
  \prim{O} = \text{octahedron}
\]

\subsection{Boolean operations}
\label{subsec:ops}

\[
  \sub = \text{smooth-subtract} \qquad
  \add = \text{smooth-union}    \qquad
  \intr = \text{smooth-intersect}
\]

All three operations accept a blending radius $\rho \geq 0$; at $\rho = 0$ they
reduce to standard sharp Boolean operations.

% ══ 3 ══════════════════════════════════════════════════════════════════════════
\section{Step Notation}
\label{sec:step}

\subsection{Parameters}
\label{subsec:params}

Each iteration step carries three real parameters
$(\sigma, \delta, \rho)$ defined in \autoref{tab:params}.

\begin{table}[h]
\centering
\caption{Step parameters.}
\label{tab:params}
\begin{tabular}{cllr}
\toprule
Symbol & JSON field & Meaning & Omit when \\
\midrule
$\sigma$ & \code{scale\_factor}    & Child size relative to parent   & never \\
$\delta$ & \code{distance\_factor} & Placement along symmetry axis   & $\delta = 1$ \\
$\rho$   & \code{smooth\_radius}   & Boolean blending width          & $\rho = 0$ \\
\bottomrule
\end{tabular}
\end{table}

\subsection{Step formula}
\label{subsec:step_formula}

\[
  \boxed{\; \mathrm{op}_{\prim{P}}^{\,\sigma}(\delta)_{\rho} \;}
\]

\noindent Examples:
\[
  \sub_{\prim{S}}^{0.5}{}_{\!0.02}
  \qquad
  \add_{\prim{C}}^{0.3}(1.2)_{\!0.05}
  \qquad
  \intr_{\prim{O}}^{0.45}
\]

% ══ 4 ══════════════════════════════════════════════════════════════════════════
\section{Full Grammar Formula}
\label{sec:formula}

\[
  \boxed{\;
    \Gamma \;=\; \G{G}[\prim{P_0}] \;:\; s_1 \cdot s_2 \cdots s_n
  \;}
\]

where $\G{G}$ is the symmetry group (\autoref{subsec:groups}),
$\prim{P_0}$ the seed primitive (\autoref{subsec:primitives}),
and $s_1, \ldots, s_n$ are steps (\autoref{sec:step}) applied in order.
Steps act on \emph{new} nodes by default; a trailing ${}^{\!\star}$
marks \code{apply\_to: all}.

% ══ 5 ══════════════════════════════════════════════════════════════════════════
\section{Named Grammars}
\label{sec:named}

\paragraph{\code{sierpinski\_classic}}
Classic tetrahedral Sierpi\'nski gasket — sphere seed, alternating
subtract/union/subtract with decreasing smoothness.
\[
  \Gamma_{\mathrm{classic}}
  = \G{T_d}[\prim{S}]
  : \sub_{\prim{S}}^{\half}{}_{\!0.02}
  \cdot \add_{\prim{S}}^{\half}{}_{\!0.01}
  \cdot \sub_{\prim{S}}^{\half}{}_{\!0.005}
\]

\paragraph{\code{sierpinski\_octahedron}}
Same alternating pattern, octahedron seed and primitives under $\G{T_d}$.
\[
  \Gamma_{\mathrm{oct}}
  = \G{T_d}[\prim{O}]
  : \sub_{\prim{O}}^{\half}{}_{\!0.02}
  \cdot \add_{\prim{O}}^{\half}{}_{\!0.01}
  \cdot \sub_{\prim{O}}^{\half}{}_{\!0.005}
\]

\paragraph{\code{sierpinski\_cube}}
Subtractive cube fractal under $\G{T_d}$.
\[
  \Gamma_{\mathrm{cube}}
  = \G{T_d}[\prim{C}]
  : \sub_{\prim{C}}^{\half}{}_{\!0.02}
  \cdot \add_{\prim{C}}^{\half}{}_{\!0.01}
  \cdot \sub_{\prim{C}}^{\half}{}_{\!0.005}
\]

\paragraph{\code{sierpinski\_void}}
Four-step icosahedral grammar producing high-genus void structure.
\[
  \Gamma_{\mathrm{void}}
  = \G{I_h}[\prim{S}]
  : \sub_{\prim{S}}^{\half}{}_{\!0.02}
  \cdot \add_{\prim{S}}^{\half}{}_{\!0.01}
  \cdot \sub_{\prim{S}}^{\half}{}_{\!0.005}
  \cdot \add_{\prim{S}}^{\half}{}_{\!0.01}
\]

% ══ 6 ══════════════════════════════════════════════════════════════════════════
\section{Mutation and Crossover Shorthand}
\label{sec:mutation}

\subsection{Point mutation}
\label{subsec:point_mutation}

A grammar reporting a delta from a named parent:
\[
  \Gamma' = \Gamma_{\mathrm{classic}}
  \;\Big[\;
    s_1 \mapsto \sub_{\prim{S}}^{0.47}{}_{\!0.025},\quad
    s_3 \mapsto \sub_{\prim{C}}^{0.53}{}_{\!0.003}
  \;\Big]
\]

\subsection{Crossover}
\label{subsec:crossover}

A crossover product with single cut point at position 2:
\[
  \Gamma' = \Gamma_{\mathrm{classic}}[{:}2] \mathbin{|} \Gamma_{\mathrm{void}}[2{:}]
\]
meaning: steps 1--2 from $\Gamma_{\mathrm{classic}}$, steps 3--4 from
$\Gamma_{\mathrm{void}}$.

% ══ 7 ══════════════════════════════════════════════════════════════════════════
\section{Fitness Annotation}
\label{sec:fitness}

A fully evaluated grammar is annotated with its fitness vector:
\[
  \Gamma_{\mathrm{void}} \xrightarrow{\;\mathcal{F}\;}
  \bigl(\,f = 0.734,\; d_f = 2.51,\; t_{\min} = 0.82\,\mathrm{mm}\bigr)
\]
where $f$ is total fitness (see \href{journey.pdf}{\emph{journey.tex}} for
the full weight table), $d_f$ the estimated fractal dimension, and $t_{\min}$
the minimum wall thickness at $80\,\mathrm{mm}$ scale.

% ══ 8 ══════════════════════════════════════════════════════════════════════════
\section{POSIX Filename Slugs}
\label{sec:slugs}

The display notation uses \code{+}, \code{:}, \code{@} which are
shell-unfriendly. Gallery filenames use a safe bijective encoding
(\autoref{tab:slugmap}):

\begin{table}[h]
\centering
\caption{Display $\to$ slug character map.}
\label{tab:slugmap}
\begin{tabular}{ccl}
\toprule
Display & Slug & Meaning \\
\midrule
\code{-} & \code{N} & subtract (Nix) \\
\code{+} & \code{P} & add (Plus) \\
\code{x} & \code{X} & intersect \\
\code{:47} & \code{k47} & scale $\sigma = 0.47$ \\
\code{@1.1} & \code{d110} & distance $\delta = 1.10$ \\
\bottomrule
\end{tabular}
\end{table}

\noindent Example:

\medskip
\begin{tabular}{ll}
Display: & \code{Td.S. -s4 +s2 -s1} \\
Slug:    & \code{Td.S.Ns4Ps2Ns1} \\
Filename: & \code{rank\_01\_Td.S.Ns4Ps2Ns1.glb} \\
\end{tabular}

\end{document}
